\documentclass[11pt]{article}
\usepackage[T1]{fontenc}
\usepackage[utf8]{inputenc}
\usepackage{lmodern}
\usepackage[margin=1in]{geometry}
\usepackage{amsmath,amssymb,amsthm,mathtools,mathrsfs}
\usepackage{booktabs,longtable,array,microtype,enumitem}
\usepackage[hidelinks]{hyperref}
\hypersetup{pdftitle={Fauzan's determinant construction for zeta(5): formalization working edition},
 pdfauthor={Original mathematics: Aabir Fauzan; working edition prepared for Christopher D. Long},
 pdfsubject={Expanded mathematical proof draft, exact certificates, and autonomous Lean specification},
 pdfkeywords={zeta(5), irrationality, Hankel determinant, Lean, formalization}}
\usepackage{fancyhdr}
\setlength{\headheight}{14pt}
\pagestyle{fancy}\fancyhf{}
\fancyhead[L]{\small Fauzan's $\zeta(5)$ construction}
\fancyhead[R]{\small Formalization working edition}
\fancyfoot[C]{\thepage}
\setlength{\parskip}{4pt}
\setlength{\parindent}{1.25em}
\setcounter{tocdepth}{2}
\emergencystretch=2em
\newtheorem{theorem}{Theorem}[section]
\newtheorem{proposition}[theorem]{Proposition}
\newtheorem{lemma}[theorem]{Lemma}
\newtheorem{corollary}[theorem]{Corollary}
\theoremstyle{definition}
\newtheorem{definition}[theorem]{Definition}
\newtheorem{specification}[theorem]{Specification}
\theoremstyle{remark}\newtheorem{remark}[theorem]{Remark}
\newcommand{\Z}{\mathbb Z}\newcommand{\Q}{\mathbb Q}
\newcommand{\R}{\mathbb R}\newcommand{\C}{\mathbb C}\newcommand{\N}{\mathbb N}
\newcommand{\Qp}{\mathbb Q_p}\newcommand{\Zp}{\mathbb Z_p}
\newcommand{\vG}{v_p^{\mathrm G}}
\newcommand{\pp}[1]{\left(#1\right)_+}
\newcommand{\floor}[1]{\left\lfloor #1\right\rfloor}
\newcommand{\fracpart}[1]{\left\{#1\right\}}
\newcommand{\Hfive}[1]{H^{(5)}_{#1}}
\newcommand{\tailD}{D_{\mathrm{tail}}}
\newcommand{\taus}{\tau^{\mathrm{ext}}}
\newcommand{\tauan}{\tau^{\mathrm{an}}}
\newcommand{\gid}[1]{\par\noindent\textit{Stable target: }\path{#1}.\par\smallskip}
\newcommand{\source}[1]{\textit{Source: #1.}}
\DeclareMathOperator{\rank}{rank}
\DeclareMathOperator{\cont}{cont}
\DeclareMathOperator{\diag}{diag}
\DeclareMathOperator{\supp}{supp}
\DeclareMathOperator{\lcm}{lcm}
\numberwithin{equation}{section}

\title{Fauzan's determinant construction for $\zeta(5)$\\[3pt]
\large Expanded proof draft and Lean formalization specification}
\author{Original manuscript: Aabir Fauzan\\
\small Working edition prepared for Christopher D. Long}
\date{23 September 2026\\\small Source manuscript dated 17 September 2026}
\begin{document}
\maketitle
\begin{abstract}
This document reconstructs and expands the argument in Aabir Fauzan's
\emph{$\zeta(5)$ is irrational}, using the supplied 32-page PDF as the
mathematical source. It preserves the determinant, rational normalization,
comparison measure, and numerical constants. The local functional is given
an explicit domain, bounded division and compatibility are proved, the
residue-class distribution identity is expanded, and the inner determinant
weights are checked by a complete source--row case analysis. All finite data
in the potential and prime-sum certificates are supplied, together with an
independently executed exact-arithmetic verifier. Stable declarations,
external-theorem interfaces, and an autonomous formalization ledger are
included. This is a mathematical working edition, not a completed Lean
formalization or an author-approved revision. The original mathematical
claims remain attributed to Fauzan; numerical experiments reported in the
supplied audits are not used as proofs.
\end{abstract}
\tableofcontents
\newpage

\section{Purpose, conventions, and the main targets}
\label{sec:overview}

\subsection{Source and editorial scope}
The source is the uploaded file \texttt{ZETA5\_IS\_IRRATIONAL.pdf}, by Aabir
Fauzan, dated 17 September 2026, 32 pages. Its SHA-256 digest is
\begin{center}\small\ttfamily
\detokenize{a42c05ccc46b9a7c36428913453c4617202c59109713ba20f464ab9ec640d0da}.
\end{center}
The original source is cited as \cite{Fauzan}. No original publisher \LaTeX{} source is assumed. Source theorem numbers
are stated in the titles or surrounding text; the numbering in this working
edition is independent. Authorial attribution is not transferred by this
reconstruction.

Two distinctions are essential. First, the $+$ after the large parenthesis
in source (5.9) is a \emph{subscript denoting the positive part}. The
indicator which follows \emph{multiplies} that expression. The PDF does not
contain an added indicator. The earlier suggestion to delete an additive
indicator was a reading error, not an error in the source. We use an
unambiguous piecewise definition in \eqref{eq:douter}.

Second, $K\ge200M^2$ is an \emph{admissibility} condition for the explicit
normalization, not an effective threshold for its asymptotic decay. For
$M=200$ it says $n\ge200000$ when $K=40n$; the decay assertions may require
a larger, unspecified threshold. For $M=100000$ the admissibility threshold
is $n\ge50000000000$.

\subsection{Basic conventions and external inputs}
Let
\[
 \xi=\sum_{m=1}^{\infty}\frac1{m^5}=\zeta(5),\qquad
 \alpha=\frac3{40},\quad \lambda=\frac{37}{40},\quad
 H=\lambda+3\alpha=\frac{23}{20}.
\]
The series definition is a convenient real-valued formalization target.
The identification with the real value of the analytically continued zeta
function must be proved if the latter is used for the public endpoint.
Bernoulli numbers have convention $B_1=-1/2$; the standard identities used below are collected in \cite{DLMF}. An empty product is $1$, an
empty sum is $0$, and $u_+=\max(u,0)$ for real $u$.

All additive valuations are normalized by $v_p(p)=1$ and $v_p(0)=+\infty$.
For $A\in\Qp[X]$, $v_p^{\mathrm G}(A)$ is the minimum coefficient
valuation, with the same convention at zero. Half-integer weight
inequalities can be represented in Lean by multiplying every inequality
by two, rather than by changing the value group.

The following standard inputs are used with their exact conventional
hypotheses: Bernoulli reflection, difference, and multiplication identities;
von Staudt--Clausen; Euler's even-zeta formula; Hermite's Hurwitz-zeta
integral; the prime number theorem in the form
\begin{equation}\label{eq:pnt}
 \vartheta(T):=\sum_{p\le T}\log p,\qquad \vartheta(T)/T\longrightarrow1;
\end{equation}
and standard finite-dimensional linear algebra, complete valued-field
analysis, and measure/integration results. Their needed interfaces are
listed in Appendix~\ref{app:lean}. No assertion is made here that a
particular current Mathlib declaration already supplies any interface.
The Gaussian energy identity and determinant integral used below are
proved here rather than treated as unnamed external results.

\begin{theorem}[Main polynomial target; source Theorem 2.1]
\label{thm:polynomialtarget}
For integers $M\ge40$ and $K\in40\Z_{>0}$ satisfying $K\ge200M^2$, the
explicit polynomial $Q_{K,M}$ defined in Sections~\ref{sec:construction}
and~\ref{sec:normalization} belongs to $\Z[X]$, has degree
$h=37K/40$, and satisfies $Q_{K,M}(\xi)>0$. There exist integers
$n_{200}\ge200000$ and $n_{100000}\ge50000000000$ such that
\begin{align}
 0<Q_{40n,200}(\xi)&<\exp(-139n^2/5) &&(n\ge n_{200}),\label{eq:decay200}\\
 0<Q_{40n,100000}(\xi)&<\exp(-7907n^2/100)
 &&(n\ge n_{100000}).\label{eq:decaylarge}
\end{align}
\gid{Zeta5.polynomial_estimates}
\end{theorem}

\begin{lemma}[Integer-polynomial irrationality criterion]
\label{lem:criterion}
Let $x\in\R$, $d,c>0$, and let $P_n\in\Z[X]$ for all sufficiently large
positive integers $n$. Suppose $\deg P_n\le dn$ and
$0<P_n(x)<\exp(-cn^2)$. Then $x$ is irrational.
\gid{Zeta5.irrational_of_integer_polynomials}
\end{lemma}
\begin{proof}
If $x=a/b$ with $a\in\Z$ and $b\in\Z_{>0}$, let $D_n$ be an integer
degree bound with $D_n\le dn$. The evaluation identity
\[
 b^{D_n}P_n(a/b)=\sum_{j=0}^{D_n}[X^j]P_n\,a^j b^{D_n-j}
\]
shows this is an integer. It is positive and smaller than
$\exp(dn\log b-cn^2)$, which tends to zero. Eventually it is a positive
integer less than one, a contradiction. Equivalently use
$D_n=\deg P_n$, since positivity excludes the zero polynomial.
\end{proof}

\begin{theorem}[Fauzan's irrationality target; source Theorem 1.1]
\label{thm:main}
The real number $\zeta(5)$ is irrational.
\gid{Zeta5.zeta_five_irrational}
\end{theorem}
\begin{proof}
Apply Theorem~\ref{thm:polynomialtarget} and
Lemma~\ref{lem:criterion} with $d=37$ and $c=139/5$.
The proof of the polynomial target is completed in Section~\ref{sec:finish}.
\end{proof}

\begin{corollary}[Optional approximation target; source Corollary 1.2]
\label{cor:approx}
There is $b_0\in\Z_{>0}$ such that
$|\zeta(5)-a/b|>b^{-260}$ for all $a\in\Z$ and all integers $b\ge b_0$.
\gid{Zeta5.rational_approximation_260}
\end{corollary}
This target, proved in Appendix~\ref{app:approx}, is independent of the
formalization of the main irrationality endpoint: it must not delay that
endpoint or be silently included among its assumptions.

\section{The functional, determinant, degree, and positivity}
\label{sec:construction}
Throughout, choose $n\in\Z_{>0}$ and put
\begin{equation}\label{eq:parameters}
 K=40n,\qquad N=3n,\qquad h=37n=K-N.
\end{equation}
Define
\[
 D_m(t)=\prod_{j=1}^m(t+j^2),\quad D_0=1,\qquad
 \Hfive j=\sum_{v=1}^j v^{-5},\quad \Hfive0=0.
\]
Let $\mathscr R_-$ be the $\Q$-vector space of rational functions in $t$
whose poles, if any, are simple and lie among $-j^2$, $j\ge1$. A function
in this space has a unique polynomial part and a unique finite simple
partial-fraction expansion.

\begin{definition}[The rational functional]
\label{def:mu}
The $\Q$-linear map $\mu_X:\mathscr R_-\to\Q[X]$ is defined by
\begin{align}
 \mu_X(t^e)&=\frac{(-1)^e B_{2e+2}(2e+3)(2e+4)(2e+5)}{24}
 &&(e\ge0),\label{eq:mumom}\\
 \mu_X\left(\frac1{t+j^2}\right)
 &=j^4(X-\Hfive j)-\frac14+\frac1{2j}
 &&(j\ge1).\label{eq:mupole}
\end{align}
Its restriction to polynomials is denoted $\mu$.
\gid{Zeta5.muX}
\end{definition}
The output is affine in $X$. The domain is a vector space, not a ring
closed under unrestricted multiplication: repeated poles are not silently
allowed.

Set
\begin{align}
 G_K(X)&=\left[\mu_X\left(\frac{D_N(t)^6t^{i+j}}{D_K(t)}\right)
             \right]_{0\le i,j<h},&
 \Delta_K(X)&=\det G_K(X),\label{eq:G}\\
 S_K&=\frac{(K!)^{2h}4^{h-1}}
 {(N!)^{12h}\prod_{i=1}^{h-1}((2i)!)^2},&F_K&=S_K\Delta_K.
 \label{eq:scalar}
\end{align}
After cancelling $D_N$, the rational input is $W(t)t^{i+j}/\tailD(t)$,
where $W=D_N^5$ and $\tailD=D_K/D_N$. Its $h$ possible poles are distinct,
so Definition~\ref{def:mu} applies. Section~\ref{sec:normalization} defines
a positive rational $m_{K,M}$ and
\begin{equation}\label{eq:Qdef}
 Q_{K,M}=m_{K,M}F_K.
\end{equation}

\begin{proposition}[Exact degree; source (2.9)]
\label{prop:degree}
The polynomial $\Delta_K$ has degree $h$, and
\begin{equation}\label{eq:leading}
 [X^h]\Delta_K=(-1)^{h(h-1)/2}
       \prod_{j=N+1}^{K}j^4D_N(-j^2)^5\ne0.
\end{equation}
Every nonzero scalar multiple, including $F_K$ and $Q_{K,M}$ when defined,
has the same degree.
\gid{Zeta5.determinant_degree}
\end{proposition}
\begin{proof}
Write $s_j=-j^2$, $N<j\le K$, and $V_{ij}=s_j^i$. Extracting the
coefficient of $X$ from each entry gives
\[
 [X]G_K=V\diag\left(\frac{j^4D_N(s_j)^5}{\tailD'(s_j)}\right)V^{\mathsf T}.
\]
The coefficient of $X^h$ in a determinant of affine entries is the
determinant of this coefficient matrix. Since
\[
 \det(V)^2=\prod_{j<k}(s_k-s_j)^2,\qquad
 \prod_j\tailD'(s_j)=(-1)^{h(h-1)/2}\prod_{j<k}(s_k-s_j)^2,
\]
the Vandermonde factors cancel and give \eqref{eq:leading}. All remaining
factors are nonzero. The degree is at most $h$ by multilinearity, and hence
is exactly $h$.
\end{proof}

\begin{proposition}[Positive moment representation; source Proposition 2.2]
\label{prop:moments}
Put
\begin{equation}\label{eq:weight}
 w(y)=\frac{(2\pi)^4y^5}{12}\sum_{\ell=1}^{\infty}\ell^4e^{-2\pi\ell y}
 \qquad(y>0).
\end{equation}
For every $R\in\mathscr R_-$,
\begin{equation}\label{eq:momentrep}
 \mu_{\xi}(R)=\int_0^\infty R(y^2)w(y)\,dy.
\end{equation}
The integral is absolutely convergent. The matrix $G_K(\xi)$ is real
symmetric positive definite; in particular $\Delta_K(\xi)>0$.
\gid{Zeta5.positive_moment_representation}
\end{proposition}
\begin{proof}
The weight is positive. The identity
$\sum_{\ell\ge1}\ell^4q^\ell=q(1+11q+11q^2+q^3)/(1-q)^5$ shows that
$w(y)\to1/\pi$ as $y\downarrow0$ and
$w(y)=O((1+y)^5e^{-2\pi y})$ as $y\to\infty$. Thus each polynomial or
permitted pole term is absolutely integrable; the same holds for their
finite sums.

For $e\ge0$, monotone convergence and the gamma integral give
\[
 \int_0^\infty y^{2e}w(y)\,dy
 =\frac{(2e+5)!}{12(2\pi)^{2e+2}}\zeta(2e+2),
\]
which is \eqref{eq:mumom} by Euler's formula.

For $a>0$, let $f(y)=(e^{2\pi y}-1)^{-1}$. Termwise differentiation on
